\documentclass[journal]{IEEEtran}

% ===================== 宏包 =====================
\usepackage{amsmath,amssymb,amsfonts,amsthm}
\usepackage{algorithm,algorithmic}
\usepackage{graphicx}
\usepackage{booktabs,multirow,array}
\usepackage[table]{xcolor}
\usepackage{colortbl}
\usepackage{pifont}
\usepackage{cite}
\usepackage{url}
\usepackage[hidelinks]{hyperref}
\usepackage{amsmath, amssymb}

\newtheorem{definition}{Definition}
\newtheorem{proposition}{Proposition}
\newtheorem{lemma}{Lemma}
\newtheorem{assumption}{Assumption}

\definecolor{best}{HTML}{D4EDDA}
\definecolor{second}{HTML}{D1ECF1}
\definecolor{rowhi}{HTML}{FFF3CD}
\definecolor{negbad}{HTML}{F8D7DA}
\definecolor{posgood}{HTML}{D4EDDA}
\definecolor{tabgray}{HTML}{F5F5F5}

\newcommand{\cmark}{\checkmark}
\newcommand{\xmark}{$\times$}
\newcommand{\best}[1]{\cellcolor{best}{\bfseries #1}}
\newcommand{\secondb}[1]{\cellcolor{second}#1}
\newcommand{\pos}[1]{\textcolor{green!50!black}{{\bfseries #1}}}
\newcommand{\bad}[1]{\textcolor{red!70!black}{#1}}
\newcommand{\dON}[1]{\cellcolor{green!18}{\bfseries #1}}
\newcommand{\dOFF}[1]{\cellcolor{gray!22}#1}
\newcommand{\gcheck}{{\color{green!50!black}\ding{51}}}
\newcommand{\rcross}{{\color{gray!65}\ding{55}}}
\newcommand{\heat}[2]{\cellcolor{blue!#1!white}#2}
\newcommand{\ind}{\mathbb{I}}

\newcommand{\system}{SpecSched}

\begin{document}

% ===================== 文档信息 =====================
\title{SpecSched: Cache-Coupled Speculative Scheduling for Multi-Tenant LLM Serving}

\author{
\IEEEauthorblockN{\parbox[t]{0.96\linewidth}{\centering
Kangkang Sun, \IEEEmembership{Member,~IEEE}, Jianhua Li, \IEEEmembership{Senior Member,~IEEE}, Xiuzhen Chen, \IEEEmembership{Member,~IEEE}, and Minyi Guo, \IEEEmembership{Fellow,~IEEE}}}
\thanks{Kangkang Sun, Jianhua Li, Xiuzhen Chen and Minyi Guo are
with the Shanghai Key Laboratory of Integrated Administration Technologies for Information Security, School of Computer Science, Shanghai Jiao Tong University, Shanghai 200240, China (e-mail: szpsunkk@sjtu.edu.cn; lijh888@sjtu.edu.cn; chenxz@sjtu.edu.cn; guo-my@cs.sjtu.edu.cn).}
% \thanks{Junyi He is  School of Science and Engineering, The Chinese University of Hong Kong, Shenzhen, Guangdong, China (e-mail: junyihe@link.cuhk.edu.cn).}
\thanks{\textit{The corresponding author is Jianhua Li.}}

% Submission / arXiv: edit \submissionjournal; after arXiv registration, set \arxivurl and uncomment the second \thanks.
% \newcommand{\submissionjournal}{an IEEE journal}% e.g., the IEEE Transactions on Mobile Computing
% \newcommand{\arxivurl}{}% e.g., https://arxiv.org/abs/2605.XXXXX
% \thanks{This manuscript has been submitted to \submissionjournal\ for possible publication. Copyright may be transferred without notice, after which this version may no longer be accessible.}%
% \thanks{A preprint version is available at \url{\arxivurl}.}%
}

\maketitle

% ============================================================
%                       ABSTRACT
% ============================================================
\begin{abstract}
Speculative decoding lets a small draft model propose tokens that a larger target model verifies in parallel. In shared serving, draft length also changes verifier load, padding, queueing, and key-value cache residency, so locally profitable drafts can reduce service-level-objective goodput. We treat resident key-value state as both memory occupancy and reusable progress: eviction can convert accepted-token gain into re-prefill work, verifier congestion, and remote-cache movement. A real edge-cloud diagnostic exposes this coupling: speculation reduces target invocations by 57.8\% and raises verification yield from 1.00 to 2.37 tokens per target call, yet uncoordinated remote drafting lowers throughput to 0.18 times cloud-only decoding. We formulate multi-tenant speculative serving as online goodput maximization under verifier and cache constraints. \system{} uses resident-aware admission, marginal verifier allocation, padding-aware batching, future-value eviction, and warm-up fallback. On ShareGPT/LMSYS trace replay with Qwen2.5-7B/0.5B, \system{} tracks but does not statistically exceed the best fixed-draft envelope. Its measured strength is cache-visible control: it raises goodput by 69.5\% to 93.5\% over NoSpec and by 2.6\% to 19.0\% over common-executor TETRIS-Proxy, EAGLE-2-Proxy, and SpecServe-Proxy scheduler-signal baselines, not native-system reproductions. NoEvict loses 68.7\% high-load goodput. Measured non-stationary transfer replay and a calibrated route/KV simulation bound the remaining scope.
\end{abstract}

\begin{IEEEkeywords}
Speculative decoding, large language models, multi-tenant serving, 
KV-cache, resource management.
\end{IEEEkeywords}

% ============================================================
%                    1. INTRODUCTION
% ============================================================
\section{Introduction}
\label{sec:intro}

Large language models (LLMs) are increasingly deployed as shared online
services that serve many concurrent requests under latency targets.
Speculative decoding~\cite{leviathan2023fast, chen2023accelerating}
reduces autoregressive target calls by letting a lightweight draft model
propose candidate tokens that the target model verifies in parallel.
This token-level mechanism is effective when a request is considered in
isolation, but multi-tenant serving introduces a hidden
cache--computation coupling: the draft length chosen for one request
also changes verifier load, batch padding, queueing delay, and
KV-cache residency for other requests. As a result, a system can accept
more draft tokens yet deliver lower end-to-end goodput when shared
KV-cache pressure or verifier contention converts the saved target
calls into waiting and re-prefill work.

Recent speculative decoding methods improve the draft itself: tree
search expands candidate paths~\cite{miao2024specinfer}, feature-level
heads raise acceptance~\cite{li2024eagle, li2024eagle2}, and multi-head
decoders add parallel proposals~\cite{cai2024medusa}. These mechanisms
do not by themselves determine whether a shared serving system meets
latency targets. Under multi-user contention, their benefit depends on
draft length, verification batching, and KV-cache management together:
longer drafts consume verifier compute and memory, batch composition
changes throughput and per-user latency, and evicting a resident
context invalidates draft progress and triggers costly re-prefill.

Serving-level speculative systems start to manage these interactions.
TETRIS~\cite{wu2025tetris} selects draft tokens under a fixed
verification budget, FASER~\cite{chen2026faser} manages speculative
length under dynamic workloads, and DiP-SD~\cite{xu2026dipsd} studies
distributed edge drafting with a shared verifier. Their scheduling
logic still leaves KV-cache residency mostly inside the serving engine.
That separation is costly. Keeping a request's cache determines its
memory footprint, future prefill cost, ability to continue drafting,
and contribution to verifier congestion. A scheduler that maximizes
immediate accepted tokens can still reduce long-term goodput if it
causes cache thrashing or repeated re-prefill. The control loop spans
epochs: draft choices change future KV pressure, eviction changes
future draft value, and SLO urgency can amplify re-prefill work when
the cache is already saturated.

This paper addresses the remaining serving-level challenge: how to
coordinate speculation when each draft decision changes both the current
verification batch and the future cache state. No existing system
jointly optimizes speculation depth, verification batching, and
KV-cache eviction under one online control loop. The missing cache
signal affects three operational choices: whether the deployed
target/draft pair is profitable, how deeply each request should draft
under a shared verifier, and which resident contexts should survive
when eviction changes future re-prefill cost and speculative value.
The technical difficulty is that these effects are measured on different
timescales: accepted tokens are observed in the current verification
epoch, while eviction and re-prefill change whether later epochs can
continue drafting at all. Treating cache management as a post-hoc memory
policy therefore loses the quantity that matters to speculative serving,
namely the future accepted-token value protected by each resident KV
block.
This view changes the cache signal available to the scheduler: a
resident KV block should be valued by the speculative progress it
preserves, not only by recency, size, or current memory pressure.
The paper studies this idea on one verifier because that is the
measured artifact, but the abstraction is intentionally resource-centric:
the same resident state is the local instance of a broader KV resource
graph whose value can also depend on routing likelihood, remote reuse,
and migration cost in disaggregated serving.

The central contribution of this paper is a cache-coupled serving
controller that turns this valuation into an executable scheduling
policy. \system{} jointly selects candidate requests, assigns marginal
draft tokens, builds padding-aware verification batches, and evicts
cached states by their future speculation value per KV block. These
components use the same currency, expected SLO-weighted accepted tokens
per unit of verifier and cache pressure, so the scheduler can trade off
an extra speculative token, a padded verification slot, and a re-prefill
event inside one online control loop. A throughput-ratio detector
separately disables speculation when a target/draft pair is empirically
unprofitable during warm-up.

We make this contribution concrete through three supporting results:

\begin{itemize}
    \item \textit{Coupling diagnosis and formulation.}
    We identify speculation-cache coupling in a real edge-cloud
    diagnostic: collaborative speculation reduces target invocations by
    $57.8\%$ and raises verification yield from $1.00$ to $2.37$ tokens
    per target call, yet lowers end-to-end throughput to $0.18\times$
    of cloud-only decoding when overheads are uncontrolled. We then
    formulate multi-tenant speculative serving as online SLO-aware
    goodput maximization under verifier-compute, routing, and KV-cache
    constraints. The measured artifact fixes routing to one verifier,
    while the formulation keeps the route as an explicit control.

    \item \textit{Cache-residency-aware control.}
    We express the online decision with four controls:
    route $r$, draft length $k$, verification batch $\mathcal{B}$, and
    eviction set $\mathcal{E}$. The implemented single-verifier system
    uses the local route and optimizes the remaining three controls.
    The resulting scheduler combines resident-aware admission, marginal
    verifier allocation, padding-aware batching, and eviction by future
    speculation value per KV block. Each module instantiates the same
    cache-aware speculation principle rather than an independent
    heuristic.

    \item \textit{Measured benefits and scope boundaries.}
    We implement \system{} on vLLM and evaluate it on
    ShareGPT/LMSYS-derived trace replay with calibrated heterogeneous
    SLOs and explicit KV-cache sweeps.
    \system{} remains statistically tied with the offline-tuned
    BestFixed envelope, while raising SLO-aware goodput over NoSpec by
    $69.5\%$ to $93.5\%$ for $N=16$ to 64 across three seeds, adding
    $2.6\%$ to $10.0\%$ over TETRIS-Proxy, $14.1\%$ to $19.0\%$ over
    EAGLE-2-Proxy, and $12.8\%$ to $14.1\%$ over SpecServe-Proxy. The
    evaluation identifies when online control reduces offline
    draft-length tuning, when homogeneous traces make fixed drafts
    sufficient, and when cache-aware admission is needed to avoid
    re-prefill cascades. The proxy comparisons are common-executor
    scheduler-signal studies rather than native-system reproductions.
\end{itemize}

% ============================================================
%              2. RELATED WORK
% ============================================================
\section{Related Work}
\label{sec:related-work}

Related work spans speculative decoding, LLM serving and KV-cache
management, and speculative serving systems. The closest speculative
systems adapt draft length or batch speculative work. The closest
serving stacks expose KV-cache residency as a core runtime resource.
\system{} targets the missing intersection: online draft scheduling
whose utility depends on which KV states remain resident under SLO
constraints.

\textit{Speculative decoding and collaborative drafting.}
Speculative decoding was introduced by Leviathan et
al.~\cite{leviathan2023fast} and Chen et
al.~\cite{chen2023accelerating}, who showed that a lightweight draft
model can propose multiple tokens and a larger target model can verify
them in one forward pass without changing the output distribution.
Following this principle, later work improves draft quality or
parallelism through feature-level heads~\cite{li2024eagle,li2024eagle2},
multi-head drafting~\cite{cai2024medusa}, tree search
~\cite{miao2024specinfer,chen2024sequoia}, staged pipelines
~\cite{spector2023accelerating}, online adaptation~\cite{liu2024online},
and big-little decoder pairs~\cite{kim2023speculative}. These methods
increase token-level yield, but they mainly target a single request or
local verification batch and do not model shared-verifier contention or
KV-cache growth across users.

\textit{LLM serving and KV-cache management.}
LLM serving systems optimize batching, memory management, and
distributed execution. Orca~\cite{yu2022orca} introduced
iteration-level batching, vLLM~\cite{kwon2023efficient} proposed
PagedAttention, SGLang exposes radix-tree cache reuse for structured
program execution~\cite{zheng2024sglang}, and recent systems
disaggregate prefill/decode~\cite{zhong2024distserve,patel2024splitwise},
schedule chunked prefill~\cite{agrawal2024taming}, migrate
requests~\cite{sun2024llumnix}, or improve
fairness~\cite{sheng2024fairness}. KV-cache work reduces memory
pressure through heavy-hitter eviction~\cite{zhang2023h2o},
disaggregation~\cite{qin2024mooncake}, compression
~\cite{liu2024cachegen,feng2025evicpress}, reuse
~\cite{gim2024prompt,cheng2025lmcache}, offloading
~\cite{xiong2024layerkv}, and admission or concurrency signals
~\cite{jaillet2025kvsched,chen2026concur}. These systems show why
KV-cache residency should be a first-class serving resource. Their
control objectives, however, are not tied to speculative draft value:
they do not jointly choose per-request draft lengths, verification
batches, and KV eviction by the expected future speculation value of
resident states.

\textit{Parallel and distributed scheduling view.}
The same gap can be stated as an online multi-resource scheduling
problem. Disaggregated and dynamic serving systems allocate prefill,
decode, migration, and batching work across shared GPU pools
~\cite{zhong2024distserve,patel2024splitwise,agrawal2024taming,
sun2024llumnix}, while fairness and cache-constrained schedulers study
latency objectives under queueing, memory, and concurrency pressure
~\cite{sheng2024fairness,jaillet2025kvsched,chen2026concur}. In
speculative serving, the resource vector additionally includes draft
latency, verifier padding, temporary speculative KV, permanent resident
KV, and eviction-induced re-prefill. \system{} contributes at this
systems layer: it exposes KV residency to the speculative scheduler so
that draft-depth and cache-admission decisions are made in the same
online control loop.

\textit{Draft adaptation for speculative serving.}
Recent systems move speculative decoding from a single-request primitive
to a serving-time decision. TETRIS~\cite{wu2025tetris} selects
promising draft tokens under a fixed verification budget,
FASER~\cite{chen2026faser} manages phase overlap and speculative
length, DiP-SD~\cite{xu2026dipsd} studies edge drafting with a shared
verifier, TurboSpec~\cite{liu2024turbospec} adapts draft length to
batch size and memory headroom, and EXSPEC~\cite{zhang2025batchspec}
formalizes correctness concerns in batched speculation. These systems
show that serving context matters for speculation. Their primary
control variables are draft tokens, phase overlap, batching feasibility,
or output equivalence. \system{} instead asks when a cached state is
worth keeping because it preserves future speculative progress and
avoids re-prefill under SLO pressure.

\textit{SLO-aware speculative serving and edge interference.}
SpecServe~\cite{huang2025specserve} adapts speculation under SLOs,
WISP~\cite{li2026wisp} suppresses edge-side waste and interference, and
FairBatching~\cite{lyu2025fairbatching} studies service fairness in
continuous batching. These works connect speculation or batching to
service objectives, but they do not make eviction and re-prefill part of
the online speculation objective. \system{} is therefore complementary,
not a native replacement for these systems: it treats KV-cache residency
and eviction as speculation-control decisions, not only as
memory-management side effects.

Table~\ref{tab:comparison} summarizes this positioning.

\begin{table}[t]
\centering
\caption{Positioning against closely related speculative-serving
systems. Proxy rows are common-executor scheduler-signal comparisons,
not native-system reproductions.}
\label{tab:comparison}
\scriptsize
\setlength{\tabcolsep}{2.0pt}
\renewcommand{\arraystretch}{1.12}
\resizebox{\columnwidth}{!}{%
\begin{tabular}{@{}p{0.22\columnwidth}p{0.31\columnwidth}ccc@{}}
\toprule
\multirow{2}{*}{System} & \multirow{2}{*}{Speculative-serving control} &
\multicolumn{3}{c}{Cache-coupled objective} \\
\cmidrule(lr){3-5}
& & SLO & KV residency & Re-prefill / evict \\
\midrule
TETRIS~\cite{wu2025tetris}
  & batch draft-token selection & \rcross & \rcross & \rcross \\
TurboSpec~\cite{liu2024turbospec}
  & closed-loop draft control & goodput & partial & \rcross \\
FASER~\cite{chen2026faser}
  & phase and draft-length control & partial & \rcross & \rcross \\
DiP-SD~\cite{xu2026dipsd}
  & distributed edge drafting & \rcross & batch feas. & \rcross \\
WISP~\cite{li2026wisp}
  & edge dynamic drafting and batching & \gcheck & \rcross & \rcross \\
SpecServe~\cite{huang2025specserve}
  & adaptive SLO-aware speculation & \gcheck & \rcross & \rcross \\
FairBatching~\cite{lyu2025fairbatching}
  & fairness-aware batch formation & \gcheck & inherit & inherit \\
vLLM~\cite{kwon2023efficient}
  & continuous batching and KV blocks & \rcross & \gcheck & spec.-obliv. \\
\midrule
\rowcolor{rowhi}{\bfseries \system{} (ours)}
  & admission, $k$, batch, and eviction & \gcheck & \gcheck & \gcheck \\
\bottomrule
\end{tabular}
}
\end{table}

\textit{Positioning of this paper.}
\system{} focuses on the cache side of serving-level speculation.
Ignoring residency can trigger re-prefill, evict high-yield requests,
or inflate padding. Unlike systems that only tune draft depth or batch
tokens, \system{} makes cache residency part of the online speculation
decision. The evaluation therefore asks when this added control is
worthwhile, and when simpler fixed-depth or eviction policies are
sufficient.
This distinction is stricter than appending a cache heuristic after draft
selection. In \system{}, residency is an input to candidate admission,
marginal draft allocation, batch shaping, and eviction, because the
same cached state determines both current memory feasibility and future
speculation option value. The paper therefore compares scheduler
signals under a common executor and reports native integration only for
the vLLM prototype.


\section{Motivating Benchmark}
\label{sec:motivating-benchmark}

We use a real prototype benchmark to show why token-level speculation
gains do not automatically become serving gains. The cloud uses
Qwen2.5-7B-Instruct on an RTX~3090~Ti as the target model, while
one Jetson device runs Qwen2.5-0.5B as the draft worker. We
generate $N=\{4,8,16,32\}$ concurrent client streams, compare
cloud-only decoding with fixed-$k$ edge-cloud speculation at
$k=3$, and use ShareGPT traces with a 128-token prompt cap and 8
generated tokens.

\begin{figure}[t]
    \centering
    \includegraphics[width=\columnwidth]{figures/figure1_motivation_singlecol.pdf}
    \caption{Motivating prototype benchmark with a single 
    Jetson edge draft worker and one cloud target server.}
    \label{fig:motivation}
\end{figure}

\textit{Observation 1: target-call savings do not guarantee
serving throughput.}
Figure~\ref{fig:motivation}(a) shows that cloud-only decoding
maintains about 46~tok/s, while fixed-$k$ edge-cloud speculation
reaches only 8.5 to 8.6~tok/s. At $N=32$, throughput drops from
46.57 to 8.54~tok/s, or about $0.18\times$ of cloud-only serving.

Queueing and tail latency explain this loss.
Figures~\ref{fig:motivation}(b) and~\ref{fig:motivation}(c)
show that edge-cloud speculation has $5.43\times$ to
$6.44\times$ higher P99 TPOT and faster-growing P95 queueing
delay. The bottleneck is not draft acceptance but the combined
service time of edge drafting, communication, serialized
verification, and queueing.

The collaborative path still has token-level value.
Figure~\ref{fig:motivation}(d) shows that edge-cloud speculation
raises verification yield from 1.00 to 2.37 tokens per target
call. The draft acceptance rate is about 50\%, and target calls per
generated token fall from 1.000 to 0.422. The Jetson draft model
therefore provides useful candidates, but remote drafting,
communication, and queueing consume the saved target calls.

\textit{Implication.}
The right abstraction is not \emph{speculation versus no speculation},
but online control of speculation under shared serving resources.
A scheduler should decide when collaboration is useful, how many
draft tokens to generate, which requests to verify together, and
which KV-cache states to keep resident.

% ============================================================
%                 4. BACKGROUND
% ============================================================
\section{Background}
\label{sec:background}

The formulation uses two ingredients: speculative decoding and
multi-tenant LLM serving. Once requests share verifier compute and
KV-cache memory, token-level speculation decisions become scheduling
decisions.

\subsection{Speculative Decoding}
\label{sec:bg-specdecode}

Speculative decoding~\cite{leviathan2023fast, 
chen2023accelerating} proceeds in rounds. In each round:

\begin{enumerate}
    \item The draft model $M_d$ autoregressively 
    generates $k$ candidate tokens 
    $\hat{y}_1, \ldots, \hat{y}_k$.
    
    \item The target model $M_t$ verifies all 
    $k$ candidates in a single forward pass.
    
    \item A rejection sampling procedure accepts a 
    prefix of length $a \leq k$, where 
    $\Pr[a \geq j]$ depends on the distributional 
    agreement between $M_d$ and $M_t$.
\end{enumerate}

The expected number of accepted tokens per round is:
\begin{equation}
    \label{eq:expected-accepted}
    \mathbb{E}[a] = \sum_{j=1}^{k} 
    \prod_{\ell=1}^{j} \alpha_\ell,
\end{equation}
where $\alpha_\ell$ is the acceptance probability at 
position $\ell$. If $\alpha_\ell = \alpha$ for all $\ell$ (i.i.d.\ 
approximation), then:
\begin{equation}
    \label{eq:iid-accepted}
    \mathbb{E}[a] = \frac{1 - \alpha^{k+1}}{1 - \alpha} 
    - 1.
\end{equation}

Speculative speedup depends on the draft-to-target cost ratio and
the acceptance profile. Acceptance is user-dependent and
position-dependent, and later draft positions often become less
reliable as prediction errors compound. In multi-tenant settings, this
heterogeneity makes a single fixed draft length suboptimal for at least
some users.

\subsection{Multi-Tenant LLM Serving}
\label{sec:bg-serving}

Production LLM serving systems such as 
vLLM~\cite{kwon2023efficient} and 
Orca~\cite{yu2022orca} employ continuous batching 
to serve multiple concurrent requests. Key resources are as follows:

\begin{itemize}
    \item \textit{Target model compute:} bounded by GPU 
    FLOPs. Verification of $B$ requests with draft 
    lengths $k_1, \ldots, k_B$ requires processing on 
    the order of $\sum_{i=1}^B (k_i{+}1)$ token 
    positions (draft verification plus bonus logits).
    
    \item \textit{KV-cache memory:} bounded by GPU 
    high-bandwidth memory (HBM). Following vLLM's 
    PagedAttention~\cite{kwon2023efficient}, each request's 
    context is stored in fixed-size blocks of 
    $B_{\text{blk}}$ tokens. The number of blocks for user 
    $i$ with context length $L_i$ is 
    $b_i(L_i) = \lceil L_i / B_{\text{blk}} \rceil$.
\end{itemize}

\begin{figure}[t]
    \centering
    \includegraphics[width=\columnwidth]{figures/background2.pdf}
    \caption{Overview of multi-tenant speculative 
    decoding. $N$ users share one target model and 
    bounded KV-cache. The scheduler jointly determines 
    draft lengths $\{k_i\}$, verification batch $B$, 
    and eviction decisions $\{e_i\}$ at each epoch.}
    \label{fig:system-overview}
\end{figure}

% ============================================================
%              5. SYSTEM MODEL & PROBLEM FORMULATION
% ============================================================
\section{System Model and Problem Formulation}
\label{sec:system-model}

We formalize only the state needed for online decisions. At each epoch,
the general control plane chooses a route $r$, draft lengths $k$, a
verification batch $\mathcal{B}$, and an eviction set $\mathcal{E}$.
The route selects where verification and KV reuse occur. In the
measured vLLM prototype, $r$ is fixed to the local verifier, reducing
the executable decision to $k$, $\mathcal{B}$, and $\mathcal{E}$. These
controls must satisfy compute, routing, and KV-cache constraints while
maximizing SLO-aware goodput. Even the simplified single-epoch
admission-and-draft decision is NP-hard.

\subsection{System Model}
\label{sec:model-system}

Concurrent users share verifier capacity and bounded KV-cache
resources. Each user's runtime state records only the quantities used
by the scheduler: context length, an acceptance estimate, SLO urgency,
cache residency, and the serving route. We consider a serving system
with:

\begin{itemize}
    \item A set of candidate routes $\mathcal{R}$. A route can denote a
    local verifier, a verifier pool, or a remote KV path. The current
    artifact instantiates the degenerate route set
    $\mathcal{R}=\{r_{\mathrm{local}}\}$.

    \item A target model $M_t$ with route-dependent verifier token
    budget $C_v^{r}(t)$ per scheduling epoch (maximum draft tokens
    verified in parallel on route $r$).
    
    \item A KV-cache of $M_{\text{KV}}^r$ blocks managed by 
    vLLM's block allocator~\cite{kwon2023efficient}, with 
    block size $B_{\text{blk}}$ tokens per block. For the measured
    single-verifier artifact, $M_{\text{KV}}^r=M_{\text{KV}}$.
    
    \item A set of active users 
    $\mathcal{N}(t) \subseteq [N]$ at time $t$, 
    each with an ongoing request.
\end{itemize}

This is written as a route-aware model, but the measured system uses
only the local route. In a disaggregated serving stack, $r$ would encode
a node or pool index, a remote-KV hit path, and a fetch or migration
cost. The present prototype does not optimize those cluster-level
choices. Instead, it isolates the common control signal that such a
scheduler would still need: whether preserving a resident KV state is
worth the compute, memory, and future re-prefill cost it avoids.

\begin{definition}[User State]
At epoch $t$, each active user $i \in \mathcal{N}(t)$ is 
characterized by the state vector:
\begin{equation}
    s_i(t) = \bigl(L_i(t), \,
    \hat{\alpha}_{i}(t), \,
    Q_i(t), \,
    h_i(t), \,
    r_i(t)\bigr),
\end{equation}
where $L_i(t)$ is the current context length, 
$\hat{\alpha}_{i}(t)$ is the estimated draft-token acceptance
probability, $Q_i(t)$ is an SLO urgency virtual queue, and 
$h_i(t) \in \{0,1\}$ indicates whether user $i$'s 
KV-cache is resident ($h_i=1$) or evicted ($h_i=0$). The route
$r_i(t)\in\mathcal{R}$ records where request $i$ is served. In the
implementation evaluated in this paper, $r_i(t)=r_{\mathrm{local}}$
for all requests.
\end{definition}

Let $a_i(k_i,t)$ denote the number of \emph{accepted draft 
tokens} in epoch $t$. The number of 
\emph{generated tokens} per speculation round is
\begin{equation}
    \label{eq:generated-tokens}
    \tilde{a}_i(k_i,t) = 1 + a_i(k_i,t),
\end{equation}
where the \emph{$+1$} is the target-model bonus token produced 
after rejection sampling. The expected generated tokens are:
\begin{equation}
    \label{eq:Gi}
    G_i(k_i, t) = \mathbb{E}[\tilde{a}_i(k_i,t)] 
    = 1 + \sum_{j=1}^{k_i} 
    \hat{\alpha}_{i}(t)^j.
\end{equation}

This scalar estimate is the control signal used by the current
scheduler. A position-level estimator can be plugged into the same
interface, but the main text keeps the model minimal.
% \textit{SLO virtual queue.}
To incorporate heterogeneous SLOs online, each user 
maintains:
\begin{equation}
    \label{eq:virtual-queue}
    Q_i(t+1) = \min\!\left(
    \Bigl[ Q_i(t) + 
    \text{TPOT}_i(t) - \tau_i^{\text{SLO}} \Bigr]^+,
    Q_{\max}\right),
\end{equation}
where $\tau_i^{\text{SLO}}$ is user $i$'s TPOT target. 
For users not verified in epoch $t$, $\text{TPOT}_i(t)$ is 
replaced by an estimated delay deficit 
$\Delta t/\widehat{G}_i(t)$, where 
$\widehat{G}_i(t)$ is the expected generated-token yield if 
the user were served in the current epoch. The cap $Q_{\max}$ 
prevents requests that are already far beyond their SLO from 
monopolizing the scheduler and repeatedly inducing re-prefill. 
The base SLO weight is $w_i(t) = 1 + \rho \, Q_i(t)$ with 
urgency parameter $\rho > 0$.

Because an SLO-critical request whose KV-cache has been evicted 
can consume extra prefill work before making progress, we 
use a cache-resident-aware scheduling weight:
\begin{equation}
    \label{eq:resident-weight}
    \bar{w}_i(t) =
    w_i(t)\,[\epsilon_h+(1-\epsilon_h)h_i(t)],
\end{equation}
where $\epsilon_h \in (0,1)$ gives evicted requests a nonzero 
but discounted priority. This discount is applied consistently 
in candidate selection and draft allocation. In addition, each 
epoch has a re-prefill admission budget:
\begin{equation}
    \label{eq:reprefill-budget}
    R_{\text{prefill}}(t) =
    \max\!\left(\left\lfloor r_{\max}\, C_v(t) \right\rfloor, 1\right),
\end{equation}
where $r_{\max}$ is the maximum fraction of verifier slots that 
may be consumed by requests whose KV-cache is currently evicted. 
This guardrail reduces the positive feedback loop in which SLO 
violations raise priority, high-priority evicted requests trigger 
re-prefill, and re-prefill creates more queueing delay.

\begin{definition}[Scheduling Decision]
At each scheduling epoch $t$, the scheduler determines:
\begin{enumerate}
    \item \textit{Route:}
    $\mathbf{r}(t)=(r_1(t),\ldots,r_N(t))$, where
    $r_i(t)\in\mathcal{R}$ selects the verifier and KV path for
    request $i$. The current artifact fixes $r_i=r_{\mathrm{local}}$.

    \item \textit{Draft lengths:} 
    $\mathbf{k}(t) = (k_1(t), \ldots, k_N(t))$, 
    where $k_i \in \{0, 1, \ldots, k_{\max}\}$. 
    Setting $k_i = 0$ means user $i$ is paused.
    
    \item \textit{Verification batch:} 
    $\mathcal{B}(t) \subseteq \mathcal{N}(t)$, the 
    subset of users whose drafts are verified in 
    this epoch.
    
    \item \textit{Eviction set:} 
    $\mathcal{E}(t) \subseteq \mathcal{N}(t)$, 
    users whose KV-cache is evicted to free memory.
\end{enumerate}
\end{definition}

\subsection{Performance Metrics}
\label{sec:model-metrics}

We use SLO-aware goodput as the primary objective because it 
captures both throughput and latency satisfaction. Unlike raw 
throughput, SLO-aware goodput only counts generated tokens from 
requests whose TPOT satisfies their latency target.

% \textit{Per-user goodput.}
For user $i$ verified at epoch $t$, let 
$T_{\text{draft},i}(k_i)$ denote the draft generation 
latency for $k_i$ tokens, $T_{\text{verify}}(\mathcal{B}(t))$ 
the shared verification latency for the entire batch 
$\mathcal{B}(t)$, and $T_{\text{wait},i}(t)$ the queuing 
delay before verification. The per-user goodput 
(generated tokens per unit wall-clock time) is:
\begin{equation}
    \label{eq:per-user-goodput}
    g_i(t) = \frac{\tilde{a}_i\bigl(k_i(t)\bigr)}
    {T_{\text{draft},i}(k_i) + 
     T_{\text{verify}}(\mathcal{B}(t)) + 
     T_{\text{wait},i}(t)}.
\end{equation}

% \textit{System goodput.}
Let $\Delta t$ denote the wall-clock duration of epoch $t$. 
The SLO-aware system goodput is:
\begin{equation}
    \label{eq:sys-goodput}
    \mathrm{Goodput} = 
    \frac{\sum_t \sum_{i \in \mathcal{B}(t)} 
    \ind\bigl[\text{TPOT}_i(t) \leq \tau_i^{\text{SLO}}\bigr]
    \cdot \tilde{a}_i\bigl(k_i(t)\bigr)}
    {\sum_t \Delta t}.
\end{equation}

% \textit{SLO attainment.}
User $i$ meets its SLO if its time-per-output-token 
(TPOT) satisfies:
\begin{equation}
    \label{eq:slo}
    \text{TPOT}_i(t) = 
    \frac{T_{\text{draft},i} + T_{\text{verify}} + 
    T_{\text{wait},i}}{\mathbb{E}[\tilde{a}_i(k_i,t)]} 
    \leq \tau_i^{\text{SLO}},
\end{equation}
where the denominator uses $G_i(k_i,t)$ in scheduling 
and $\tilde{a}_i(k_i,t)$ in realized measurements. 
This avoids division by zero when all draft tokens are 
rejected but the target model still emits the bonus token.

The system-wide SLO attainment rate is:
\begin{equation}
    \label{eq:slo-attainment}
    S = \frac{1}{|\mathcal{N}|} \sum_{i \in \mathcal{N}} 
    \ind\bigl[\text{TPOT}_i \leq \text{SLO}_i\bigr].
\end{equation}

\subsection{Resource Constraints}
\label{sec:model-constraints}

The scheduler must satisfy four coupled constraints in every 
epoch: route capacity, verifier-token capacity, block-level KV-cache
capacity, and padding-aware verification latency. These constraints
determine whether locally beneficial draft tokens can actually 
improve system-level goodput.

\textit{Route capacity.}
Let $\mathcal{B}_r(t)=\{i\in\mathcal{B}(t):r_i(t)=r\}$ be the requests
assigned to route $r$. A route may represent a local verifier in the
current prototype or a verifier/KV path in a distributed deployment.
For each route:
\begin{equation}
    \label{eq:route-capacity}
    \mathcal{B}(t) = \bigcup_{r\in\mathcal{R}}\mathcal{B}_r(t),
    \qquad
    \mathcal{B}_r(t)\cap\mathcal{B}_{r'}(t)=\emptyset
    \;\; (r\neq r').
\end{equation}
The measured artifact uses a single route, so this constraint reduces
to $\mathcal{B}_{r_{\mathrm{local}}}(t)=\mathcal{B}(t)$.

% \textit{Verifier token budget.}
The target model verifies $k_i$ draft positions and 
produces one bonus-token logit per batched user, so the 
effective verification load is closer to $k_i{+}1$ than 
$k_i$. For each route $r$:
\begin{equation}
    \label{eq:compute-constraint}
    \sum_{i \in \mathcal{B}_r(t)} \bigl(k_i(t)+1\bigr)
    \leq C_v^r(t).
\end{equation}

When profiling is available, this can be replaced by 
$\sum_{i \in \mathcal{B}_r(t)}
c_v^r\bigl(k_i(t),L_i(t)\bigr) \leq C_v^r(t)$.

\textit{Block-level KV-cache constraint.}
Following vLLM's PagedAttention, the number of KV blocks 
occupied by user $i$ is:
\begin{equation}
    \label{eq:block-count}
    b_i(L_i) = \left\lceil \frac{L_i}{B_{\text{blk}}} 
    \right\rceil.
\end{equation}

We separate permanent residency from transient 
speculative memory:
\begin{equation}
    \label{eq:kv-decomp}
    M_{\text{used}}^r(t) = M_{\text{perm}}^r(t) + 
    M_{\text{temp}}^r(t) + M_{\text{act}}^r(t),
\end{equation}
where
\begin{align}
    M_{\text{perm}}^r(t) &= 
    \sum_{i: h_i(t)=1,\,r_i(t)=r} b_i\bigl(L_i(t)\bigr), \\
    M_{\text{temp}}^r(t) &= 
    \sum_{i \in \mathcal{B}_r(t)} b_i^{\text{tmp}}\bigl(k_i(t)\bigr), \\
    M_{\text{act}}^r(t) &\text{ captures route-local activation / scratch buffers.}
\end{align}
Before verification,
\begin{equation}
    \label{eq:memory-constraint}
    M_{\text{used}}^r(t) \leq M_{\text{KV}}^r,
    \qquad \forall r\in\mathcal{R}.
\end{equation}
After verification, permanent context grows by generated 
tokens:
\begin{equation}
    \label{eq:context-update}
    L_i(t{+}1) = L_i(t) + \tilde{a}_i\bigl(k_i(t)\bigr),
\end{equation}
while rejected draft suffixes release the corresponding
$M_{\text{temp}}^r$.

\textit{Padding-aware verification cost.}
Variable-length drafts in a batch incur padding overhead. 
We model the padding cost of batch $\mathcal{B}$ as:
\begin{equation}
    \label{eq:pad-cost}
    \begin{aligned}
    \text{PadCost}(\mathcal{B}) 
    &= |\mathcal{B}| \max_{i \in \mathcal{B}} k_i 
    - \sum_{i \in \mathcal{B}} k_i \\
    &\quad + \zeta \left( 
    |\mathcal{B}| \max_{i \in \mathcal{B}} L_i 
    - \sum_{i \in \mathcal{B}} L_i \right),
    \end{aligned}
\end{equation}
where $\zeta \geq 0$ weights context-length imbalance. 
Verification latency $V(\mathcal{B}, \mathbf{k})$ grows 
with $\text{PadCost}(\mathcal{B})$ and batch size.

% \textit{Eviction score.}
If memory is over budget, \system{} ranks resident contexts by the
future speculative progress preserved by each KV block:
\begin{equation}
    \label{eq:eviction-cost}
    \text{KeepValue}_i(t) =
    \frac{\bar{w}_i(t)\,G_i(k_i^\star,t)}
    {b_i(L_i)+\mu\,\ell_{r_i}(t)} ,
\end{equation}
where $k_i^\star$ is the draft length that the allocator would assign
if request $i$ remained resident, $\ell_{r_i}(t)$ is a route-cost term
for remote fetch, migration, or route switching, and $\mu\geq0$ maps
that route cost into the block-cost denominator. In the measured local
artifact, $\ell_{r_i}(t)=0$, so Eq.~\eqref{eq:eviction-cost} reduces to
expected SLO-weighted generated tokens per resident KV block.
\system{} evicts the lowest-value resident contexts until the block
budget is feasible. Re-prefill risk is handled by the resident-aware
weight and the per-epoch re-prefill admission budget rather than by a
multi-term eviction reward.
The route-cost term makes explicit how the same score would account for
local reuse, remote reuse, and migration. We keep it zero in all
reported experiments to avoid claiming unmeasured cluster-level routing
behavior.

\subsection{Joint Optimization Problem}
\label{sec:model-optimization}

We combine the performance objective and resource constraints 
into a single per-epoch scheduling problem. The formulation 
chooses routes, draft lengths, verification batches, admission 
decisions, and cache evictions jointly, making explicit the 
tradeoff between immediate accepted-token gain, routing cost, and
future cache stability.

At each epoch $t$, the scheduler chooses draft lengths 
$\mathbf{k}$, routes $\mathbf{r}$, verification batch
$\mathcal{B}$, admission set $\mathcal{A}$, and eviction set
$\mathcal{E}$. We 
formulate the per-epoch objective as:
\begin{equation}
    \label{eq:joint-opt}
    % \boxed{
    \begin{aligned}
    \max_{\mathbf{r},\mathbf{k}, \mathcal{B}, \mathcal{A}, \mathcal{E}} 
    \quad & \sum_{i \in \mathcal{B}} \bar{w}_i(t) \, G_i(k_i, t) 
    - \lambda_v V(\mathcal{B}, \mathbf{k}) 
    - \lambda_p \text{PadCost}(\mathcal{B}) \\
    & \quad - \lambda_r\sum_{i\in\mathcal{B}}\ell_{r_i}(t)
    + \lambda_e \sum_{i \notin \mathcal{E}} 
    h_i(t)\,\text{KeepValue}_i(t) \\
    \text{s.t.} \quad 
    & \mathcal{B}=\bigcup_{r\in\mathcal{R}}\mathcal{B}_r,
    \quad \mathcal{B}_r\cap\mathcal{B}_{r'}=\emptyset, \\
    & \sum_{i \in \mathcal{B}_r} \bigl(k_i{+}1\bigr) \leq C_v^r(t), 
    \quad \forall r \in \mathcal{R} \\
    & M_{\text{used}}^r(t) \leq M_{\text{KV}}^r, 
    \quad \forall r \in \mathcal{R} \\
    & \sum_{i \in \mathcal{B}} \ind[h_i(t)=0]
    \leq R_{\text{prefill}}(t), \quad \forall t \\
    & \mathbb{E}[\text{TPOT}_i(t)] \leq \tau_i^{\text{SLO}} 
    \quad \text{(via virtual queues)}, \\
    & k_i \in \{0, 1, \ldots, k_{\max}\}, 
    \quad r_i\in\mathcal{R}, \quad \forall i
    \end{aligned}
    % }
\end{equation}
where $\lambda_v$, $\lambda_p$, $\lambda_r$, and $\lambda_e$ balance 
verification latency, padding overhead, route cost, and resident cache
value. In the evaluated single-verifier prototype,
$\mathcal{R}=\{r_{\mathrm{local}}\}$ and $\ell_{r_i}=0$, so the
implemented optimizer is the local route specialization of
Problem~\eqref{eq:joint-opt}. 
The multi-epoch problem maximizes time-averaged SLO-aware 
goodput subject to the constraints above.

\begin{lemma}[Single-Epoch Hardness]
\label{prop:hardness}
Even with known acceptance estimates, no eviction, no 
padding penalty, and fixed per-request priority weights, 
the single-epoch admission-and-draft allocation subproblem 
of Problem~\eqref{eq:joint-opt} is NP-hard when each user 
has a binary choice between being paused and being served 
with a fixed draft length.
\end{lemma}

\begin{proof}
We reduce from the 0-1 Knapsack Problem. Consider an 
instance with $N$ items, where item $i$ has cost 
$c_i \in \mathbb{Z}^+$ and value $v_i \in \mathbb{Z}^+$, 
and the knapsack capacity is $W$. Without loss of 
generality, take $c_i \geq 2$ by multiplying all item 
costs and the capacity by two.

Construct a single-epoch instance of 
Problem~\eqref{eq:joint-opt} as follows. For each item 
$i$, create one request with two choices: pause it, or 
serve it with fixed draft length $\bar{k}_i=c_i-1$, so 
its verification load $\bar{k}_i+1$ equals $c_i$. Set 
$C_v=W$, set all verification-latency, padding, and 
eviction penalties to zero, and relax the SLO constraints. 
Let $\omega_i=v_i/G_i(\bar{k}_i,t)$ be the fixed priority 
weight of request $i$, which is well-defined because 
$G_i(\bar{k}_i,t)>0$. If priorities must take the specific 
virtual-queue form $w_i(t)=1+\rho Q_i(t)$, multiply all item 
values by a common positive constant so that every 
$\omega_i\geq 1$. This scaling does not change the 
knapsack optimizer, and the initial virtual queues can then 
realize the constructed priorities.

The single-epoch objective reduces to:
\begin{equation}
  \begin{aligned}
  \max_{\{x_i\}} \quad & 
  \sum_{i=1}^{N} x_i \omega_i G_i(\bar{k}_i,t) \\
  \text{s.t.} \quad & 
  \sum_{i=1}^{N} x_i(\bar{k}_i{+}1) \leq C_v, \\
  & x_i \in \{0,1\}.
  \end{aligned}
\end{equation}
which is exactly the 0-1 Knapsack Problem because 
$\omega_i G_i(\bar{k}_i,t)=v_i$ and 
$\bar{k}_i+1=c_i$. Since 0-1 Knapsack is NP-hard, the 
single-epoch admission-and-draft allocation problem is 
NP-hard.
\end{proof}

Intuitively, Lemma~\ref{prop:hardness} implies 
that no polynomial-time algorithm can find the globally 
optimal schedule in general, \emph{even for a single 
scheduling epoch without eviction}. The multi-epoch 
problem with eviction and SLO constraints is strictly 
harder. This motivates practical online scheduling rules, 
with formal statements limited to simplified subproblems 
(Section~\ref{sec:algorithm}).


% ============================================================
%              6. COUPLING ANALYSIS
% ============================================================
\section{Analysis of the Speculation-Cache Coupling Effect}
\label{sec:analysis}

Independent per-user speculation can be suboptimal when users share
verifier compute and KV-cache memory. The single-user rule recalled in
the supplementary appendix chooses $k$ by maximizing expected
accepted tokens per draft and verification cost, assuming a private
target verifier. Shared serving violates that assumption: the effective
verification cost for one request depends on the total batch load,
padding length, queueing delay, and other users' draft choices. When
many requests simultaneously choose long drafts, the verifier must
either pad a large batch or defer part of the work. Queueing delay can
then push requests past their TPOT targets, so tokens that improve raw
throughput no longer count toward SLO-aware goodput. A useful scheduler
therefore needs a load-aware draft cap, not only a per-request
acceptance estimate.

KV-cache pressure creates the second coupling. Evicting a request frees
memory immediately, but the next time that request is served it must
re-prefill before speculation can continue. Under high load this can
produce a feedback loop: SLO-urgent evicted requests are prioritized,
re-prefill consumes verifier time, queues grow, and more requests become
SLO urgent. The scheduler should therefore keep a cached state when its
future accepted-token value per KV block is high, and evict it when the
same memory can protect more useful speculative progress elsewhere.
Section~\ref{sec:impl-eval} tests this intuition directly through KV
pressure sweeps and ablations rather than relying on a closed-form
queueing model.


% ============================================================
%              7. ALGORITHM DESIGN
% ============================================================
\section{The \system{} Scheduling Algorithm}
\label{sec:algorithm}

\system{} aligns four decisions often treated separately: where a
request is served, who should be verified, how many draft tokens to
generate, and which KV-cache states should remain resident. The current
artifact fixes the route to one local verifier, while keeping the route
variable in the formulation for distributed deployments. The online
policy then has a simple greedy structure. Stage~1 selects candidate
requests by SLO urgency, acceptance estimate, and cache residency.
Stage~2 allocates draft tokens one at a time to the candidate with the
largest marginal SLO-weighted yield per compute and cache cost. Stage~3
evicts the resident contexts with the smallest future speculation value
per KV block. Together, these stages implement one abstraction: a
cached context is an option to produce future accepted tokens without
re-prefill. The scheduler spends verifier slots and cache blocks only
when that option value is large enough relative to its route, padding,
queueing, and memory cost.

\subsection{Overview}
\label{sec:alg-overview}

\system{} operates over repeated scheduling epochs
(Figure~\ref{fig:algorithm}). At epoch $t$, the request pool exposes
SLO slack, waiting time, acceptance estimates
$\hat{\alpha}_{i}(t)$, cache residency $h_i(t)$, and context
length $L_i(t)$. A route-selection layer would assign requests to a
verifier and KV path before candidate selection. In the measured
prototype this layer is the identity map to $r_{\mathrm{local}}$, so
\system{} first selects a candidate verification batch and then
allocates draft lengths within that batch, keeping draft allocation and
batch composition aligned. It also separates temporary speculative KV
from permanent residency: draft tokens consume transient KV during
verification, accepted tokens extend the permanent context, and
rejected suffixes release temporary KV.

Each epoch then has four phases: candidate batch selection using SLO
urgency, residency, waiting time, and the re-prefill budget, marginal
draft allocation under verifier and transient KV budgets, KV-cache
admission/eviction using the keep-value score in
Eq.~\eqref{eq:eviction-cost}, and target verification followed by state
updates.

\begin{figure}[t]
    \centering
    \includegraphics[width=\columnwidth]{figures/fig3_2.pdf}
    \caption{Overview of the \system{} scheduling loop. At 
    each scheduling epoch, \system{} jointly selects candidate 
    requests, allocates per-request draft lengths, composes a 
    padding-aware verification batch, and manages KV-cache 
    admission/eviction according to future speculation value and 
    overload-protection hooks. 
    Verification outcomes update acceptance estimates, context 
    lengths, SLO states, and KV-cache residency for the next 
    epoch.}
    \label{fig:algorithm}
\end{figure}

\subsection{Phase 1: Candidate Batch Selection}
\label{sec:alg-candidate}

We first select users likely to benefit from verification in this
epoch, avoiding draft allocation for requests that will not be
verified. Define a pre-scheduling score:
\begin{equation}
    \label{eq:pre-priority}
    \text{pre\_score}_i =
    \bar{w}_i(t)\,\hat{\alpha}_{i}(t),
\end{equation}
where $\bar{w}_i(t)$ is the cache-resident-aware SLO weight in
Eq.~\eqref{eq:resident-weight}. Users enter $\mathcal{C}(t)$ in
decreasing score order while reserving one verifier slot for the
bonus token. Evicted requests also consume the re-prefill budget
$R_{\text{prefill}}(t)$, and are deferred once that budget is
exhausted.

\begin{algorithm}[t]
\caption{Candidate Batch Selection}
\label{alg:candidate}
\begin{algorithmic}[1]
\REQUIRE Active users $\mathcal{N}$, budget $C_v$, 
         states $\{s_i(t)\}$, weights $\{\bar{w}_i(t)\}$
\ENSURE Candidate set $\mathcal{C}$
\STATE Compute $\text{pre\_score}_i$ via Eq.~\eqref{eq:pre-priority}
\STATE Sort users by $\text{pre\_score}_i$ descending
\STATE $\mathcal{C} \leftarrow \emptyset$, $R \leftarrow C_v$
\STATE $R_p \leftarrow R_{\text{prefill}}(t)$
\FOR{each user $i$ in sorted order}
    \IF{$R \geq 1$}
        \IF{$h_i(t)=0$ and $R_p = 0$}
            \STATE continue
        \ENDIF
        \STATE $\mathcal{C} \leftarrow \mathcal{C} \cup \{i\}$
        \IF{$h_i(t)=0$}
            \STATE $R_p \leftarrow R_p - 1$
        \ENDIF
        \STATE $R \leftarrow R - 1$ 
               \COMMENT{reserve bonus-token slot}
    \ENDIF
\ENDFOR
\RETURN $\mathcal{C}$
\end{algorithmic}
\end{algorithm}

\subsection{Phase 2: SLO-Aware Marginal Draft Allocation}
\label{sec:alg-draft}

Given $\mathcal{C}(t)$, assign draft lengths 
$\mathbf{k}(t) = (k_i)_{i \in \mathcal{C}}$ to maximize 
SLO-weighted expected generated tokens while respecting 
the remaining verifier budget $C_v(t) - |\mathcal{C}(t)|$ 
for draft positions.

Define the marginal score of assigning one additional 
draft token to user $i$ at epoch $t$:
\begin{equation}
    \label{eq:marginal-score}
    \Delta_i(k, t) = 
    \frac{\bar{w}_i(t)\,\hat{\alpha}_{i}(t)^{k+1}}
    {\Delta V_i(k, t) + \eta \, \Delta M_i(k, t)
    + \eta_p\,\ind[h_i(t)=0]},
\end{equation}
where $\Delta V_i(k,t)$ is the incremental verifier cost, 
$\Delta M_i(k,t)$ is the incremental KV-block pressure 
(from growing $L_i$), and $\eta \geq 0$ trades off compute 
and cache. The additional $\eta_p$ term discourages allocating 
long drafts to requests whose cache must first be reconstructed. 
The numerator is the cache-resident-aware SLO-weighted marginal 
generated-token gain. The denominator penalizes users whose 
next draft token congests verification, cache, or re-prefill.

We allocate draft tokens greedily: at each step, assign 
one token to the user with the highest $\Delta_i(k_i, t)$ 
until the verifier budget is exhausted or marginal gain 
falls below $\epsilon_{\min}$. Since 
$\hat{\alpha}_{i}^{k+1}$ decreases with $k$ for
$0\leq\hat{\alpha}_{i}\leq1$, this naturally balances draft
lengths across users and reduces padding variance in verification
batches.

\begin{algorithm}[t]
\caption{SLO-Aware Marginal Draft Allocation}
\label{alg:water-filling}
\begin{algorithmic}[1]
\REQUIRE Candidate users $\mathcal{C}$, budget $C_v$, 
         states $\{s_i(t)\}$, weights $\{\bar{w}_i(t)\}$
\ENSURE Draft lengths $\{k_i\}_{i \in \mathcal{C}}$
\STATE Initialize $k_i \leftarrow 0$ for all 
       $i \in \mathcal{C}$
\STATE $R \leftarrow C_v - |\mathcal{C}|$ 
       \COMMENT{draft slots after bonus reservation}
\WHILE{$R > 0$ and $\exists i \in \mathcal{C}: k_i < k_{\max}$}
    \STATE $i^* \leftarrow \arg\max_{i \in \mathcal{C}: k_i < k_{\max}} \; 
           \Delta_i(k_i, t)$ via Eq.~\eqref{eq:marginal-score}
    \IF{$\Delta_{i^*}(k_{i^*}, t) < \epsilon_{\min}$}
        \STATE break
    \ENDIF
    \STATE $k_{i^*} \leftarrow k_{i^*} + 1$
    \STATE $R \leftarrow R - 1$
\ENDWHILE
\RETURN $\{k_i\}$
\end{algorithmic}
\end{algorithm}

The scalar EMA makes the token-yield component diminish with draft
depth: the next token contributes $\hat{\alpha}_i^{k+1}$ expected
accepted tokens. The implemented scheduler uses the density score in
Eq.~\eqref{eq:marginal-score}, which also includes incremental
verification cost and KV pressure. Under realistic GPU execution,
padding makes verification latency depend on the maximum draft length
in a batch as well as context-length imbalance. The full scheduler is
therefore evaluated empirically under the measured latency path.

Since true acceptance probabilities are unknown a priori, 
we maintain an exponential moving average:
\begin{equation}
    \label{eq:alpha-estimation}
    \hat{\alpha}_{i}(t) =
    \gamma\,\frac{a_i(t-1)}{\max\{k_i(t-1),1\}}
    +(1-\gamma)\hat{\alpha}_{i}(t-1),
\end{equation}
where $\gamma \in (0,1)$ is the smoothing parameter. Position-level
feedback can replace this scalar estimator without changing the
policy interface.
This lightweight estimator tracks recent acceptance behavior for
online scheduling. Non-stationary workloads can still bias
$\Delta_i(k,t)$ rankings, so we audit the estimator separately from
the main throughput claims. As a targeted fixed-seed diagnostic, a
position-level \textit{\system{}-PosAlpha} variant changes goodput by
$-1.35\%$ to $+0.03\%$ across $N\in\{4,8,16,32,64\}$. This suggests
that, on the measured trace, load and cache coupling explain more of
the observed performance than fine-grained acceptance mis-ranking.
The EMA is therefore a lightweight control signal, not an optimal
predictor. Broader Bayes and non-stationary estimator audits remain
future work.

\subsection{Executed Verification Batch}
\label{sec:alg-batch}

After Phase~2, the executed batch is
$\mathcal{B}(t)=\{i \in \mathcal{C}(t): k_i(t)>0\}$. Unlike
standard continuous batching~\cite{yu2022orca}, speculative
verification batches have heterogeneous token counts, creating GPU
padding overhead.

The purpose of this phase is to convert per-request draft decisions
into one executable verifier call without letting a long draft from one
request inflate the latency of the whole batch. This is the point where
token-level speculation becomes a serving-level scheduling decision:
the same expected accepted tokens can have different system value
depending on the padded length and slack of the requests batched
together.

\textit{Verification latency model.}
Under the padded-and-masked implementation 
(Section~\ref{sec:varlen-verification}), the verification 
latency for batch $\mathcal{B}$ is:
\begin{equation}
  \label{eq:verify-latency}
  T_{\textup{verify}}(\mathcal{B}) = 
  f_{\textup{attn}}\!\left(|\mathcal{B}|, \, 
  k^{\mathcal{B}}_{\max}, \, 
  \bar{L}(\mathcal{B})\right) + 
  f_{\textup{ffn}}\!\left(|\mathcal{B}| \cdot 
  k^{\mathcal{B}}_{\max}\right),
\end{equation}
where $k^{\mathcal{B}}_{\max} = \max_{i \in 
\mathcal{B}} k_i$ is the padded draft length, 
$\bar{L}(\mathcal{B}) = \frac{1}{|\mathcal{B}|} 
\sum_{i \in \mathcal{B}} L_i$ is the mean context 
length, $f_{\textup{attn}}$ captures attention 
computation (quadratic in context length, linear 
in batch size), and $f_{\textup{ffn}}$ captures 
feed-forward computation (linear in total padded 
tokens). We profile $T_{\textup{verify}}$ offline for representative
$(|\mathcal{B}|, k^{\mathcal{B}}_{\max}, \bar{L})$ configurations.
In the current prototype, this profile validates
Eq.~\eqref{eq:verify-latency} as a design check for padding-aware
batching, while the online scheduler uses the measured batch path and
does not read fitted coefficients. The main experiments therefore
report measured TPOT and goodput after executing each policy, not
model-predicted service time. A cached-verification fit over 64
$(|\mathcal{B}|,k^{\mathcal{B}}_{\max},\bar{L})$ points on the 7B
target has 0.34~ms MAE, 0.42~ms RMSE, 1.42\% MAPE, and $R^2=0.975$.
An offline $\pm20\%$ perturbation audit changes at most 1.57\% of
nearby fitted-latency orderings. Because the scheduler does not use
these fitted coefficients online, the perturbation result is an
offline ranking-stability check rather than a measured goodput
robustness claim.

We formulate batch composition as a variant of 
the weighted job scheduling problem, explicitly 
incorporating padding cost (Eq.~\eqref{eq:pad-cost}):

\begin{equation}
    \label{eq:batch-opt}
    \begin{aligned}
    \max_{\mathcal{B} \subseteq \mathcal{N}} \quad 
    & \sum_{i \in \mathcal{B}} \bar{w}_i(t) \, G_i(k_i, t) 
    - \lambda_p \text{PadCost}(\mathcal{B}) \\
    \text{s.t.} \quad 
    & \sum_{i \in \mathcal{B}} (k_i{+}1) \leq C_v \\
    & T_{\textup{verify}}(\mathcal{B}) \leq 
    \min_{i \in \mathcal{B}} S_i(t), \\
    & S_i(t)=\tau_i^{\text{SLO}} G_i(k_i, t) 
    - T_{\textup{draft},i}(k_i) 
    - T_{\textup{wait},i}(t).
    \end{aligned}
\end{equation}

The SLO constraint follows by rearranging
$\textup{TPOT}_i \leq \tau_i^{\text{SLO}}$. If realized
verification latency exceeds slack, the next epoch reduces $k_i$ or
removes $i$ from $\mathcal{C}(t)$. For padding-aware tie-breaking,
we use
\begin{equation}
    \label{eq:priority}
    \text{score}_i = 
    \frac{\bar{w}_i(t) \, G_i(k_i, t)}{k_i + \epsilon},
\end{equation}
where $\bar{w}_i(t)$ encodes SLO urgency and cache residency. Users
with similar $(k_i,L_i)$ are batched together to reduce
$\text{PadCost}(\mathcal{B})$.

This design gives the scheduler a measurable advantage over
draft-length-only control. It preserves the correctness of standard
speculative verification while exposing the padding cost that determines
whether accepted tokens translate into SLO-aware goodput. As a result,
batch formation can reject a locally attractive draft when its padded
verifier cost would consume another request's SLO slack.

\subsection{Phase 3: Speculation-Aware KV-Cache Management}
\label{sec:alg-cache}

When $M_{\text{used}}(t)$ in Eq.~\eqref{eq:memory-constraint} 
would be violated (including route-local $M_{\text{temp}}^r$ during 
verification), select users for eviction that minimize the 
impact on system goodput.

The motivation for speculation-aware eviction is that a cached context
has two values at once: it avoids re-prefill and it preserves the future
ability to speculate on that request. Unlike standard LRU/LFU eviction,
we rank caches by the expected speculative progress protected per KV
block. This keeps the eviction rule simple and directly tied to the
serving objective.

We evict users with the smallest keep value:
\begin{equation}
    \label{eq:eviction-score}
    \text{EvictScore}_i = 
    \text{KeepValue}_i(t),
\end{equation}
using Eq.~\eqref{eq:eviction-cost}. This evicts caches with the 
lowest expected future generated-token value per block, not merely the
oldest cache.

The technical advantage is that eviction is expressed in the same
currency as draft allocation: expected future SLO-weighted accepted
tokens per unit of scarce verifier/cache resource. This common scoring
surface lets the scheduler trade a cache block, a re-prefill event, and
an extra speculative token inside one control loop instead of optimizing
them in separate engine components.

When score gaps are small, simple rules such as LRU or
longest-context-first can be as good or better, as the KV sweep in
Section~\ref{sec:cache-ablation} shows. The score is therefore a
structured cache-aware policy, not a claim of superiority in every
memory regime.

\begin{algorithm}[t]
\caption{Speculation-Aware Cache Eviction}
\label{alg:eviction}
\begin{algorithmic}[1]
\REQUIRE Resident users $\{i : h_i = 1\}$, 
         block budget $M_{\text{KV}}$, usage $U$ (blocks)
\ENSURE Eviction set $\mathcal{E}$
\STATE $\mathcal{E} \leftarrow \emptyset$
\WHILE{$U > M_{\text{KV}}$}
    \STATE $i^* \leftarrow \arg\min_{i : h_i = 1} 
           \text{EvictScore}_i$
    \STATE $\mathcal{E} \leftarrow 
           \mathcal{E} \cup \{i^*\}$
    \STATE $h_{i^*} \leftarrow 0$
    \STATE $U \leftarrow U - b_{i^*}(L_{i^*})$
\ENDWHILE
\RETURN $\mathcal{E}$
\end{algorithmic}
\end{algorithm}

Algorithm~\ref{alg:joint} combines the route-specialized control
decisions within each scheduling epoch: route selection, candidate
selection, marginal draft allocation, and speculation-aware cache
management. The measured implementation sets all routes to
$r_{\mathrm{local}}$. The resulting batch is then drafted, verified,
and used to update the permanent context state.

\begin{algorithm}[t]
\caption{\system{}: Joint Scheduling per Epoch}
\label{alg:joint}
\begin{algorithmic}[1]
\REQUIRE System state at epoch $t$
\ENSURE Scheduling decisions 
        $(\mathbf{r}, \mathbf{k}, \mathcal{B}, \mathcal{E})$
\STATE Set $r_i \leftarrow r_{\mathrm{local}}$ for all $i$
       \COMMENT{measured artifact, distributed deployments optimize $\mathbf{r}$}
\STATE Update $\hat{\alpha}_{i}(t)$ via 
       Eq.~\eqref{eq:alpha-estimation} for all $i$
\STATE Update $Q_i(t)$ via Eq.~\eqref{eq:virtual-queue}
\STATE Compute $\bar{w}_i(t)$ and $R_{\text{prefill}}(t)$ via
       Eqs.~\eqref{eq:resident-weight}--\eqref{eq:reprefill-budget}
\STATE $\mathcal{C} \leftarrow$ 
       \textsc{CandidateSelect}($\mathcal{N}$, $C_v$, 
       $\{s_i(t)\}$) 
       \COMMENT{Algorithm~\ref{alg:candidate}}
\STATE $\{k_i\}_{i \in \mathcal{C}} \leftarrow$ 
       \textsc{MarginalAlloc}($\mathcal{C}$, $C_v$, 
       $\{s_i(t)\}$) 
       \COMMENT{Algorithm~\ref{alg:water-filling}}
\STATE $\mathcal{B} \leftarrow \{i \in \mathcal{C}: k_i > 0\}$
\STATE Remove lowest-score evicted-cache users until
       $\sum_{i\in\mathcal{B}}\ind[h_i=0]\leq R_{\text{prefill}}(t)$
\STATE Estimate post-verification usage 
       $\widehat{M}_{\text{perm}}^{+}$ from 
       $\{G_i(k_i,t)\}_{i\in\mathcal{B}}$
\STATE Ensure 
       $\widehat{M}_{\text{perm}}^{+}+M_{\text{temp}}^r+M_{\text{act}}^r
       \leq M_{\text{KV}}$ by trimming $k_i$, removing candidates, 
       or evicting low-value resident users
\IF{$\widehat{M}_{\text{perm}}^{+} > M_{\text{KV}}$}
    \STATE $\mathcal{E} \leftarrow$ 
           \textsc{SpecEvict}($M_{\text{KV}}$, $\widehat{M}_{\text{perm}}^{+}$) 
           \COMMENT{Algorithm~\ref{alg:eviction}}
\ENDIF
\STATE Issue drafts with lengths $\{k_i\}_{i \in \mathcal{B}}$
\STATE Verify batch $\mathcal{B}$ on target model
\STATE Apply rejection sampling
\STATE Update $L_i$ via Eq.~\eqref{eq:context-update}
\STATE Release route-local $M_{\text{temp}}^r$
\IF{realized $M_{\text{perm}} > M_{\text{KV}}$}
    \STATE Invoke \textsc{SpecEvict} before admitting the next epoch
\ENDIF
\end{algorithmic}
\end{algorithm}

\subsection{Overload Detection and Degradation Hooks}
\label{sec:alg-guardrails}

\textit{Speculation-OK detector.}
Before serving a new target/draft pair, \system{} runs the same warm-up
calibration used for SLO thresholding and compares NoSpec throughput
with the best speculative calibration point. Let $G_0$ be the mean
calibration throughput at $k=0$, and let $G_{\mathrm{spec}}^{\max}$
be the best mean throughput over $k>0$. The safe variant enables
drafting only when
\begin{equation}
    \label{eq:spec-detector}
    \textsc{Speculate} \equiv
    \frac{G_{\mathrm{spec}}^{\max}}{G_0}
    \geq 1 + \epsilon .
\end{equation}
Otherwise, \textit{\system{}-Safe} keeps the same scheduler interface
but routes the policy to NoSpec. The implementation also logs the
i.i.d.\ break-even proxy
$1 + 1/(\hat{\alpha} k_{\max})$ from the calibration acceptance rate,
but uses the measured throughput ratio in Eq.~\eqref{eq:spec-detector}
because it directly includes draft-model latency, tokenizer/model
overheads, and verification batching on the deployed hardware.
We use $\epsilon=0.05$, requiring at least a 5\% warm-up throughput
gain before enabling speculation. The measured
$G_{\mathrm{spec}}^{\max}/G_0$ ratios are 0.69--0.71 for
Qwen2.5-1.5B/0.5B (OFF), 0.20--0.21 for SmolLM2-1.7B/135M (OFF),
1.32--1.44 for Pythia-2.8B/160M (ON), and 1.61--1.74 for the main
Qwen2.5-7B/0.5B pair (ON).

The virtual queue in Eq.~\eqref{eq:virtual-queue} is useful for
prioritizing requests that are close to violating their TPOT
targets, but under severe cache pressure it can interact badly
with eviction. If many SLO-critical requests have already lost
their KV-cache, blindly prioritizing them can create a feedback
loop: re-prefill consumes verifier time, queueing delay increases,
more requests violate SLOs, and their virtual queues grow. To
detect this mode, \system{} monitors a sliding-window re-prefill
rate:
\begin{equation}
    \label{eq:reprefill-rate}
    \rho_{\text{prefill}}(t) =
    \frac{\text{re-prefill events in window}}
    {\text{generated tokens in window}}.
\end{equation}

When $\rho_{\text{prefill}}(t)$ exceeds a threshold
$\theta_{\text{prefill}}$, the scheduler switches to a
conservative mode by reducing $r_{\max}$, increasing the
cache-residency discount, and clipping drafts to
$k_i \leq k_{\text{safe}}$. If SLO attainment remains zero for
$K$ windows, evicted-cache requests temporarily fall back to
ordinary target-model decoding until cache pressure drops. These
hooks preserve speculative-decoding correctness and prevent epochs
from being consumed by re-prefill work that cannot satisfy SLOs.
When queue depth, TPOT violations, and re-prefill rate are all small,
the SLO virtual queue can be bypassed in favor of FCFS ordering,
matching the ablation in Section~\ref{sec:cache-ablation}.

\textit{Arrival-regime signals.}
The implementation exposes the arrival process, arrival rate, active
request count, ready-pending backlog, oldest pending delay, and
pending-queue length to the policy interface. We implemented four
diagnostic open-loop controllers with arrival-pressure signals, with
details and measurements in the supplementary appendix. The validation shows
that these signals can stabilize open-loop behavior and reduce wasted
target-model work, but they do not yet establish a statistically
reliable goodput advantage over shallow fixed drafting. We therefore
treat arrival-pressure control as an exposed diagnostic interface rather
than as a headline contribution.

\subsection{Complexity and Analytical Scope}
\label{sec:alg-complexity}

\begin{proposition}[Per-Epoch Overhead]
\label{prop:complexity}
The scheduling overhead of Algorithm~\ref{alg:joint} 
is $O(N \cdot k_{\max} + N \log N)$ per epoch, 
dominated by the water-filling loop 
($O(N \cdot k_{\max})$) and the sorting step in 
batch composition ($O(N \log N)$). The re-prefill monitor 
and load-mode hooks use sliding-window counters and add only 
$O(N)$ work per epoch.
\end{proposition}

Since $k_{\max}$ is typically small (e.g., 
$k_{\max} \leq 20$) and $N$ is on the order of 
tens to hundreds, the scheduling overhead is 
negligible compared to model inference time 
(typically milliseconds vs.\ tens of milliseconds).

\textit{Analytical scope.}
The greedy allocation is exact only in the idealized unit-cost case
where each marginal draft token has the same verifier cost, padding is
absent, and cache pressure is ignored. Real serving violates all three
conditions. We therefore use the analysis only to justify the marginal
allocation structure and rely on measured trace replay, KV-pressure
sweeps, and ablations for performance claims. The EMA in
Eq.~\eqref{eq:alpha-estimation} is likewise a lightweight tracking
mechanism rather than a worst-case learning guarantee. Stronger
contextual estimators can plug into Eq.~\eqref{eq:marginal-score}
without changing the scheduler interface.


% ============================================================
%              8. IMPLEMENTATION AND EVALUATION
% ============================================================
\section{Implementation and Evaluation}
\label{sec:impl-eval}

We evaluate \system{} on a vLLM-based prototype with calibrated NoSpec
SLO targets, trace-derived output lengths, and explicit KV-cache budget
sweeps. The main comparison keeps the executor, traces, KV manager, and
SLO calibration fixed, so exposing residency to the speculative
scheduler is the tested variable. Proxy baselines isolate
scheduler-visible signals under the same executor. Boundary experiments
cover unsafe remote drafting, model pairs, fixed draft lengths, cache
budgets, scheduler overhead, open-loop arrivals, and route/KV stress as
diagnostics. The artifact is a single-verifier local-route study: route
costs are in the formulation, but production routing, native
speculative-serving integrations, multi-GPU scaling, and long-running
cost stability remain outside the measured scope.

\subsection{Implementation and Experimental Setup}
\label{sec:impl}
\label{sec:varlen-verification}
\label{sec:exp-setup}

We implement \system{} in vLLM~\cite{kwon2023efficient} by replacing
FCFS scheduling with candidate selection, marginal draft allocation,
padding-aware verification, and speculation-aware cache management.
The runtime hooks are the scheduler policy, the variable-length
verification path, and the block allocator interface. The target model,
draft model, token acceptance rule, and sampling distribution are
unchanged. Batched verification uses padding masks, per-request
position IDs, and block-diagonal causal masks. Following the correctness
risk identified by recent batched-speculation work~\cite{zhang2025batchspec},
we run an output-equivalence smoke test: in 1000 random prompts, batched
outputs match sequential single-request decoding bit-exactly.
The main statistical claims use Qwen2.5-7B-Instruct/Qwen2.5-0.5B-Instruct
on an RTX~3090~Ti and repeat $N\in\{16,32,64\}$ closed-loop and
heterogeneous Poisson load sweeps with seeds 7, 42, and 123. Full
implementation constants, logged fields, hardware variants, and scope
limits are moved to the supplementary appendix.

The prototype measures one target verifier with a bounded local KV pool,
plus a separate edge-cloud diagnostic for remote drafting overhead. A
deployment with prefill/decode disaggregation or KV-cache tiers would
add routing and migration terms to the same keep-value interface.

SLO-aware metrics are calibrated from NoSpec rather than chosen as
absolute constants. For each trace and concurrency level, we first run
NoSpec and set heterogeneous user targets as
\begin{equation}
  \tau_i^{\text{SLO}} =
  \beta_i \cdot \mathrm{median}\!\left(
  \text{TPOT}^{\text{NoSpec}}_{\mathrm{class}(i)}\right),
  \quad \beta_i \in \{1.25, 1.5, 2.0\}.
\end{equation}
The class index groups requests by prompt/output-length bucket. We
report the medium setting ($\beta_i=1.5$) unless otherwise stated and
sweep $\beta_i\in\{1.25,1.5,2.0\}$ for SLO sensitivity.

Baselines include real executor behavior (NoSpec, fixed-$k$
speculation, and LRU-style cache management), common-executor proxies
for EAGLE-2, TETRIS, and SpecServe/AdaSpec scheduling signals, and
\system{} variants that remove batching, eviction coupling, SLO urgency,
or a specific eviction rule. Proxy rows are not native
reimplementations. BestFixed is an offline-tuned fixed-$k$ envelope and
is treated as a boundary, not a weak baseline.
For multi-seed cells, lifts are computed per seed against the same
trace, concurrency, KV budget, and SLO calibration, then reported as
paired means with 95\% confidence intervals where space permits. When
intervals cross zero, we describe the methods as statistically tied.

\subsection{Resource Coupling Profiling}
\label{sec:profiling}

We first profile how batch size and concurrent users 
affect speculative serving resources, following the 
methodology of FASER~\cite{chen2026faser} but adding 
KV-cache metrics. The profiling records TPOT, the fraction 
of epoch time spent in verification, KV-cache block occupancy, 
and eviction/re-prefill frequency on Qwen2.5-7B-Instruct with
Qwen2.5-0.5B-Instruct drafting in vLLM. As batch size grows,
verification becomes the dominant 
bottleneck. As $N$ grows under a fixed $M_{\text{KV}}$, 
eviction and re-prefill events increase sharply, showing that
KV-cache residency is part of the speculation-control loop.

\subsection{Real Edge-Cloud Diagnostic}
\label{sec:edge-practical}

\textit{RQ1: Is uncoordinated remote drafting safe?}
We run a real distributed-draft diagnostic across seeds 7, 42, and
123: the target model remains on the RTX~3090~Ti server, while three
Jetson Orin devices serve Qwen2.5-0.5B over Tailscale/SSH. The study
uses the measured SSH/JSONL path at $N\in\{4,8\}$ with 8 requests per
cell and 32 generated tokens. It is not an optimized RPC evaluation.
Its purpose is to identify when remote drafting must be disabled before
the controlled single-GPU scheduler comparison.

\begin{figure*}[t]
    \centering
    \begin{minipage}{0.58\textwidth}
        \centering
        \includegraphics[width=\linewidth]{figures_scientific/edge_main_bottleneck_scientific.pdf}
    \end{minipage}
    \hfill
    \begin{minipage}{0.38\textwidth}
        \centering
        \includegraphics[width=\linewidth]{figures_scientific/edge_service_breakdown_n8_scientific.pdf}
    \end{minipage}
    \caption{Real edge-cloud diagnostic over Tailscale/SSH. Left:
    three-seed mean and standard deviation for SLO-aware goodput,
    P99 TPOT, and SLO attainment. Right: service-time breakdown at
    $N=8$ from the seed-7 trace. The per-phase fractions vary by less
    than $\pm 3$ percentage points across seeds 7, 42, and 123.}
    \label{fig:edge-practical}
\end{figure*}

Figure~\ref{fig:edge-practical} shows that remote drafting and RPC
dominate this unoptimized path. NoSpec reaches
$46.16\pm0.14$ and $41.88\pm7.42$~tok/s at $N=4$ and $N=8$, whereas
unprotected \system{} has zero SLO-aware goodput and P99 TPOT rises to
$643.05\pm5.94$ and $1407.31\pm6.11$~ms/token. The
Speculation-OK detector classifies the path as unsafe and
\system{}-Safe falls back to NoSpec, recovering
$45.67\pm0.42$ and $41.82\pm7.36$~tok/s with the same SLO attainment.
At $N=8$, edge drafting plus RPC accounts for about 64\% of service
time in the representative seed-7 trace. The diagnostic therefore
separates remote-path overhead from the scheduler-level gains evaluated
next.

\subsection{Main Results}
\label{sec:main-results}

\textit{RQ2--RQ3: Does cache-coupled online coordination improve
goodput, and how close is it to the offline fixed-depth envelope?}
Figure~\ref{fig:scheduler-main} reports the full load sweep for
NoSpec, the best fixed-$k$ envelope, TETRIS-Proxy, and \system{}.
The core loads $N\in\{16,32,64\}$ are repeated across three seeds. The
result supports a bounded claim: coordinated scheduling improves over
NoSpec and common-executor scheduler proxies, while remaining
statistically tied with the offline-tuned BestFixed envelope when one
draft length fits most active requests. Table~\ref{tab:baseline-performance}
lists paired goodput gains at $N=16/32/64$. Model-pair and KV-pressure
stress cells remain in the supplementary appendix.

\begin{figure*}[t]
    \centering
    \includegraphics[width=\textwidth]{figures_scientific/main_scheduler_overview_scientific.pdf}
    \caption{Scheduler evaluation under calibrated heterogeneous
    SLOs and explicit KV-cache pressure. The panels compare
    SLO-aware goodput and P99 TPOT against NoSpec, the best fixed-$k$
    envelope, TETRIS-Proxy, and \system{}.}
    \label{fig:scheduler-main}
\end{figure*}

\textit{Observation 2: \system{} improves over no-draft serving and
common-executor scheduler proxies, while matching rather than
dominating the tuned fixed-draft envelope.}
Across seeds 7, 42, and 123, \system{} improves SLO-aware
goodput over NoSpec by $69.5\%\pm4.8\%$ at $N=16$,
$93.5\%\pm6.8\%$ at $N=32$, and $91.3\%\pm6.6\%$ at
$N=64$. In absolute terms, \system{} reaches
$71.42\pm6.03$, $76.55\pm2.45$, and $76.42\pm1.80$ tok/s
at these three loads. Compared with TETRIS-Proxy, its paired
goodput lift is $2.6\%\pm0.2\%$, $5.1\%\pm1.0\%$, and
$10.0\%\pm3.1\%$, while reducing mean P99 TPOT from
585.35 to 557.27~ms/token at $N=32$ and from 1179.27 to
1028.49~ms/token at $N=64$.

Table~\ref{tab:baseline-performance} gives the main benchmark in a
compact form: evaluated baseline first, reference source second, then
paired goodput gains at $N=16/32/64$. It is not a native-system
leaderboard: the supported reading is improvement over NoSpec and
common-executor scheduler proxies, while remaining statistically tied with the offline-tuned BestFixed envelope.

% Inlined from ../paper_assets/tables/baseline_performance.tex.
\begin{table}[t]
\centering
\caption{Core baseline comparison under the common executor. Values are paired \system{} SLO-goodput lift at $N=16/32/64$ on Qwen2.5-7B/0.5B. Baseline names appear first, followed by their reference source. Proxy rows isolate scheduler-visible signals and are not native-system reproductions.}
\label{tab:baseline-performance}
\scriptsize
\setlength{\tabcolsep}{2.8pt}
\renewcommand{\arraystretch}{1.12}
\resizebox{\columnwidth}{!}{%
\begin{tabular}{@{}p{0.39\columnwidth}ccc@{}}
\toprule
\multirow{2}{*}{Baseline / comparator with source} &
\multicolumn{3}{c}{\system{} SLO-goodput lift ($\uparrow$)} \\
\cmidrule(lr){2-4}
 & $N=16$ & $N=32$ & $N=64$ \\
\midrule
NoSpec, vLLM/PagedAttention~\cite{kwon2023efficient} & \pos{+69.5\%} & \pos{+93.5\%} & \pos{+91.3\%} \\
\rowcolor{tabgray}BestFixed, Speculative decoding~\cite{leviathan2023fast,chen2023accelerating} & \secondb{$-$1.1\%} & \secondb{+0.6\%} & \secondb{+0.2\%} \\
IndepOpt, Online draft control~\cite{liu2024online,liu2024turbospec} & \pos{+28.5\%} & \pos{+30.7\%} & \pos{+18.1\%} \\
TETRIS-Proxy, TETRIS~\cite{wu2025tetris} & \pos{+2.6\%} & \pos{+5.1\%} & \pos{+10.0\%} \\
EAGLE-2-Proxy, EAGLE-2~\cite{li2024eagle2} & \pos{+14.1\%} & \pos{+19.0\%} & \pos{+15.1\%} \\
SpecServe-Proxy, SpecServe/AdaSpec~\cite{huang2025specserve} & \pos{+12.8\%} & \pos{+14.1\%} & \pos{+13.0\%} \\
\bottomrule
\end{tabular}
}
\vspace{-0.5em}
\end{table}

\textit{Common-executor proxy comparisons.}
Every proxy row uses the same target model, draft model, token
acceptance rule, verification path, KV allocator, trace replay, and
NoSpec-calibrated SLO targets. Under this controlled executor,
\system{} improves SLO-aware goodput over SpecServe-Proxy by
12.8\%, 14.1\%, and 13.0\% for $N=16,32,64$, while reducing P99 TPOT
by 9.6\%, 14.8\%, and 21.4\%. It also improves paired goodput over
EAGLE-2-Proxy by 14.1\%, 19.0\%, and 15.1\%.

\textit{Fixed-draft envelope.}
The fixed-$k$ envelope is the strongest simple comparator because it
can choose the best draft length after seeing the workload. Under
closed-loop ShareGPT/LMSYS replay, active requests often have similar
acceptance estimates, so the marginal allocator converges to the same
shallow depth selected by the offline envelope. At $N=64$, \system{} is
only $0.22\%\pm0.59\%$ above the best fixed-$k$ policy, and at $N=16$
it is $1.07\%\pm2.16\%$ below that envelope. The paired confidence
intervals cross zero. The supported claim is therefore online
adaptation that approaches the tuned envelope while avoiding manual
draft-length selection and improving over NoSpec and proxy schedulers,
not guaranteed superiority over a trace-specific fixed draft length.

\textit{BestFixed transfer diagnostic.}
Table~\ref{tab:transfer-gpu-measured} compares \system{} with the best
full-run fixed depth selected from $k\in\{1,3,5,7\}$ on three
non-stationary traces. The completed replay is mixed-to-negative: all
paired confidence intervals cross zero, and the winning fixed depth is
FixedK-1 for every seed and scenario. This indicates that the current
acceptance estimator adapts too slowly to beat a shallow full-run fixed
choice on these traces. The result is diagnostic and not headline evidence.
The calibrated stress surface is kept in the supplementary appendix.

% Inlined from ../paper_assets/tables/transfer_gpu_measured_multiseed.tex.
\begin{table}[t]
\centering
\caption{Measured three-seed non-stationary transfer replay at $N=64$. These rows are diagnostic and are not used as headline gains. BestFixed is the best full-run fixed depth from $k\in\{1,3,5,7\}$ selected per seed. Intervals are paired 95\% CIs for \system{} lift over BestFixed.}
\label{tab:transfer-gpu-measured}
\scriptsize
\setlength{\tabcolsep}{2.6pt}
\renewcommand{\arraystretch}{1.12}
\resizebox{\columnwidth}{!}{%
\begin{tabular}{@{}lcrrrr@{}}
\toprule
Scenario & BestFixed $k$ & BestFixed & \system{} & Lift & \system{} P99 \\
\midrule
Workload shift & 1/1/1 & 45.6 & 37.4 & -17.2\% [-39.6,+5.2] & 1221 \\
Bursty arrivals & 1/1/1 & 49.5 & 48.9 & -1.3\% [-4.4,+1.8] & 1213 \\
Tiered SLOs & 1/1/1 & 52.8 & 49.5 & -5.9\% [-20.3,+8.6] & 1206 \\
\bottomrule
\end{tabular}
}
\vspace{-0.5em}
\end{table}













\textit{Open-loop arrivals.}
We also replay heterogeneous Poisson arrivals at 12, 24, and
48 requests/s for $N=16,32,64$ across the same seeds. \system{}
improves goodput over NoSpec by 66.5\%, 85.4\%, and 81.1\%.
Compared with TETRIS-Proxy, paired lifts are 31.3\%, 20.9\%, and
13.6\%, with positive confidence intervals at $N=32$ and $N=64$.
Against BestFixed, the paired lifts are $-5.0\%$, $+9.0\%$, and
$+2.2\%$, all with intervals crossing zero. Open-loop pressure
therefore strengthens the case for online coordination over NoSpec
and proxy schedulers, but it still does not establish a reliable
advantage over offline-tuned fixed draft lengths.

\textit{Model-pair boundary and safe fallback.}
Additional $N=32$ model-pair and KV-pressure boundary cells are kept
outside Table~\ref{tab:baseline-performance} to avoid presenting
boundary failures as a headline ranking. The model-pair and
KV-pressure tables in the supplementary appendix show the same runtime
rule: the detector decides whether speculation should run, whereas
eviction-policy weights remain a per-regime tuning problem. On unsafe
Qwen2.5-1.5B/0.5B and SmolLM2 pairs the detector falls back near
NoSpec, while Pythia remains profitable relative to NoSpec but
statistically tied with TETRIS-Proxy and BestFixed.
Table~\ref{tab:detector-decisions} summarizes the same warm-up rule as
a deployment decision log: profitable pairs keep speculation enabled,
whereas unprofitable local and remote-draft paths fall back before they
can erase SLO-aware goodput.

% Inlined from ../paper_assets/tables/detector_decisions.tex.
\begin{table}[t]
\centering
\caption{Warm-up Speculation-OK detector decisions. The ratio column is
the measured $G_{\mathrm{spec}}^{\max}/G_0$ interval used by
Eq.~\eqref{eq:spec-detector}. The effect column reports realized
\system{}-Safe SLO-aware goodput.}
\label{tab:detector-decisions}
\scriptsize
\setlength{\tabcolsep}{3.0pt}
\renewcommand{\arraystretch}{1.12}
\begin{tabular}{@{}p{0.42\columnwidth}c c p{0.22\columnwidth}@{}}
\toprule
Target / draft pair & Ratio & Decision & Effect \\
\midrule
Qwen2.5-7B / 0.5B & 1.61--1.74 & \dON{ON} & \pos{main path enabled} \\
Pythia-2.8B / 160M & 1.32--1.44 & \dON{ON} & \pos{127.5 tok/s} \\
\midrule
Qwen2.5-1.5B / 0.5B & 0.69--0.71 & \dOFF{OFF} & \secondb{98.8 tok/s} \\
SmolLM2-1.7B / 135M & 0.20--0.21 & \dOFF{OFF} & \secondb{115.0 tok/s} \\
\rowcolor{rowhi}Jetson edge path & 0.21--0.24 & \dOFF{OFF} & \best{45.7 tok/s} \\
\bottomrule
\end{tabular}
\vspace{-0.5em}
\end{table}

\textit{Open-loop controller diagnostics.}
We implemented four diagnostic open-loop controllers with
arrival-pressure signals. They reduce wasted draft work and stabilize
backlog, but multi-seed validation does not establish a statistically
reliable goodput advantage over FixedK-1. Details and per-cell
measurements appear in the supplementary appendix.

\begin{figure}[t]
    \centering
    \includegraphics[width=0.8\columnwidth]{figures_scientific/slo_sensitivity_overview_scientific.pdf}
    \caption{SLO sensitivity under tight, medium, and loose
    TPOT targets derived from the NoSpec baseline. Each cell reports
    \system{} relative to IndepOpt for the corresponding
    $(\beta,N)$ setting.}
    \label{fig:slo-sensitivity}
\end{figure}

Calibrated SLOs preserve the benefit across tight and loose targets.
Figure~\ref{fig:slo-sensitivity} varies the SLO multiplier
$\beta$. For $N=16$, \system{} improves goodput over IndepOpt by
31.3\%, 30.0\%, and 12.3\% under tight, medium, and loose SLOs.
For $N=32$, the gains are 36.9\%, 34.0\%, and 14.0\%. At $N=64$,
the gain stays between 19.5\% and 20.2\% across the three settings.
SLO calibration changes which requests count toward goodput, while
load-aware draft caps and cache-resident admission prevent
high-concurrency re-prefill cascades.

\subsection{KV-Cache Sensitivity and Ablations}
\label{sec:cache-ablation}

\textit{RQ4: Does cache visibility change scheduling behavior under
KV pressure?}
We vary the vLLM KV budget over 256, 512, 768, and 1024 blocks and
report goodput, P99 TPOT, re-prefills, and eviction frequency in
Figure~\ref{fig:kv-sweep}. These settings exercise the cache-pressure
regime where re-prefill and eviction decisions dominate goodput.

\begin{figure*}[t]
    \centering
    \includegraphics[width=\textwidth]{figures_scientific/kv_sweep_overview_scientific.pdf}
    \caption{KV-cache budget sweep at $N=64$. The panels report
    SLO-aware goodput, P99 TPOT, and re-prefill frequency across
    eviction policies under the same trace replay and SLO calibration.}
    \label{fig:kv-sweep}
\end{figure*}

\textit{Cache-pressure diagnosis: eviction is regime dependent.}
At $N=64$, the fixed-seed sweep in Figure~\ref{fig:kv-sweep} shows
that cache-resident admission avoids the zero-goodput collapse seen in
unguided high-pressure runs. We then repeated the two representative
KV-pressure budgets, 512 and 1024 blocks, across seeds 7, 42, and
123. The three-seed results confirm that the cache policy is an
operating-regime choice rather than a single-policy win. At 512
blocks, LRU gives the highest mean goodput,
73.01$\pm$1.89~tok/s, compared with 72.10$\pm$2.91~tok/s for
\system{}. vLLM-style preemption gives the lowest P99 TPOT,
818.12$\pm$21.87~ms/token, while H2O has the fewest re-prefills,
3.18$\pm$0.29 per 1K tokens. At 1024 blocks,
longest-context-first gives the highest mean goodput,
75.02$\pm$1.78~tok/s, and the fewest re-prefills,
0.74$\pm$0.06 per 1K tokens. \system{} gives the lowest P99 TPOT at
811.34$\pm$73.06~ms/token, but its goodput is
74.28$\pm$2.13~tok/s and its re-prefill rate is
2.44$\pm$0.31 per 1K tokens. Paired confidence intervals between
\system{} and the eviction alternatives all cross zero. The implication
is narrower than a universal ranking: cache residency should be visible
to the speculation scheduler, while eviction weights should adapt to
the active memory budget and latency objective.

\begin{figure}[t]
    \centering
    \includegraphics[width=0.92\columnwidth]{figures_scientific/ablation_throughput_scientific.pdf}\\[-0.4ex]
    \centerline{\small (a) Throughput}
    \vspace{0.7ex}
    \includegraphics[width=0.92\columnwidth]{figures_scientific/ablation_p99_tpot_scientific.pdf}\\[-0.4ex]
    \centerline{\small (b) P99 TPOT}
    \caption{Ablation under high load and a constrained KV budget.
    Batched verification is the largest throughput contributor,
    while disabling eviction sharply reduces SLO-aware goodput.}
    \label{fig:ablation}
\end{figure}

Figure~\ref{fig:ablation} and Table~\ref{tab:ablation-main} test
whether the control principle is more than a draft-length heuristic.
At $N=64$, \system{} reaches 75.72~tok/s SLO-aware goodput and
916.57~ms/token P99 TPOT. Removing batched verification drops goodput
by 32.9\% and raises P99 TPOT by 30.8\%, showing that accepted tokens
must be converted into a feasible verification batch. Disabling
eviction keeps raw throughput high but reduces SLO-aware goodput by
68.7\% and SLO attainment by 0.64 because requests wait behind a
saturated cache. This is the strongest causal evidence that cache
residency cannot remain an invisible serving-engine side effect.
Disabling SLO urgency gives slightly higher goodput and lower P99 TPOT,
so load-triggered SLO urgency is preferable to always-on priority in
this regime.

% Inlined from ../paper_assets/tables/ablation_main.tex.
\begin{table}[t]
\centering
\caption{Core ablation under a constrained 384-block KV budget
($\beta=1.5$). Values are single-seed measurements from the existing
artifact. $\Delta$ columns are relative to the full \system{} row at the
same concurrency. NoSLO has slightly higher goodput than the full system
because FCFS ordering is near-optimal when acceptance rates are
homogeneous. Load-triggered SLO activation is left as future work.}
\label{tab:ablation-main}
\scriptsize
\setlength{\tabcolsep}{2.6pt}
\renewcommand{\arraystretch}{1.12}
\resizebox{\columnwidth}{!}{%
\begin{tabular}{@{}l c rr rr rr@{}}
\toprule
\multirow{2}{*}{Variant} & \multirow{2}{*}{$N$} &
\multicolumn{2}{c}{Goodput tok/s ($\uparrow$)} &
\multicolumn{2}{c}{P99 TPOT ms ($\downarrow$)} &
\multicolumn{2}{c}{SLO attain. ($\uparrow$)} \\
\cmidrule(lr){3-4}\cmidrule(lr){5-6}\cmidrule(l){7-8}
 & & Value & $\Delta$ & Value & $\Delta$ & Value & $\Delta$ \\
\midrule
\rowcolor{rowhi}\system{} (full) & 32 & {\bfseries 76.5} & -- & {\bfseries 377.5} & -- & {\bfseries 0.94} & -- \\
\system{}-NoBatch & 32 & 51.4 & \bad{$-$32.8\%} & 567.7 & \bad{+50.4\%} & 0.94 & \secondb{0.00} \\
\system{}-NoEvict & 32 & 47.3 & \bad{$-$38.1\%} & 467.0 & \bad{+23.7\%} & 0.66 & \bad{$-$0.28} \\
\system{}-NoSLO & 32 & 77.4 & \secondb{+1.1\%} & 410.5 & \bad{+8.7\%} & 0.94 & \secondb{0.00} \\
\midrule
\rowcolor{rowhi}\system{} (full) & 64 & {\bfseries 75.7} & -- & {\bfseries 916.6} & -- & {\bfseries 0.94} & -- \\
\system{}-NoBatch & 64 & 50.8 & \bad{$-$32.9\%} & 1199.2 & \bad{+30.8\%} & 0.94 & \secondb{0.00} \\
\system{}-NoEvict & 64 & 23.7 & \bad{{\bfseries $-$68.7\%}} & 893.2 & \secondb{$-$2.6\%} & 0.30 & \bad{{\bfseries $-$0.64}} \\
\system{}-NoSLO & 64 & 77.2 & \secondb{+2.0\%} & 824.9 & \secondb{$-$10.0\%} & 0.94 & \secondb{0.00} \\
\bottomrule
\end{tabular}
}
\vspace{-0.5em}
\end{table}

The ablation and pressure results separate three failure modes:
verifier saturation without batching, KV thrashing without
resident-aware control, and harmful speculation on unsafe target/draft
pairs. The supplementary appendix details the taxonomy and CPU-side
scheduler microbenchmark.

\subsection{Limitations and Scope}
\label{sec:limitations}

The evaluation supports a bounded systems claim rather than a universal
winner claim. First, BestFixed remains statistically competitive in the
measured closed-loop and heterogeneous Poisson validations: paired
confidence intervals against BestFixed cross zero at the tested loads.
The calibrated transfer diagnostic marks workloads with measured
acceptance heterogeneity, bursty arrivals, or tiered SLO mixes as
plausible fixed-depth stressors, but the three-seed transfer replay is
mixed-to-negative with intervals crossing zero. We therefore interpret
\system{} as reducing offline draft-length tuning and stabilizing cache
pressure, not as replacing a trace-specific fixed draft length in every
regime. Second, the KV sweep shows that no
static eviction rule is best across all budgets: LRU,
longest-context-first, H2O-style eviction, and vLLM-style preemption can
win different objectives. Third, all headline measurements use one 3090
Ti verifier setup, and TETRIS, EAGLE-2, and SpecServe are represented by
common-executor scheduler proxies rather than native integrations. The
supplementary appendix gives the full evidence roadmap and detailed
exclusions for production traces, multi-GPU verification, remote-KV
migration, and artifact packaging.

% ============================================================
%              9. CONCLUSION
% ============================================================
\section{Conclusion}
\label{sec:conclusion}

This paper addresses multi-tenant speculative serving as a
cache-coupled control problem in which draft length, verification
batching, and KV-cache residency must be controlled together. The main
lesson is that KV-cache residency is not only a memory-management detail:
it changes future draft value, re-prefill pressure, and SLO-aware
goodput. \system{} exposes cache residency as a first-class scheduling
variable through resident-aware admission, marginal draft allocation,
padding-aware verification, future-value eviction, and a detector that
falls back to NoSpec on unprofitable target/draft pairs.

Across three seeds, the vLLM implementation raises SLO-aware goodput
over NoSpec, TETRIS-Proxy, EAGLE-2-Proxy, and a three-load
SpecServe-Proxy sweep, while remaining statistically tied with the
offline-tuned BestFixed envelope. The transferable insight is that the
most reliable benefit of online speculative-serving control is not
dominance over a trace-specific fixed draft length, but safe adaptation:
the scheduler can approach the tuned envelope, avoid manual draft-depth
selection, and disable harmful speculation when the target/draft pair or
remote path is unprofitable.

The current evidence is bounded to the measured vLLM/3090~Ti and
Jetson settings, common-executor scheduler proxies, calibrated
route/KV simulation, and saved-output artifact checks. Next steps are
native integrations with recent speculative-serving systems, finer
component ablations of admission versus future-value eviction, stronger
batched-output equivalence audits, and a native distributed stress suite
where re-prefill is an explicit cross-pool migration cost.

% ============================================================
%              REFERENCES
% ============================================================
\bibliographystyle{IEEEtran}
\bibliography{references}
\end{document}
